\section{Limitations}

\plat{} is small and experimental. The current repository presents a compact
interpreter, examples, tests, documentation, and design notes, but it should
not be treated as a mature programming ecosystem. Its standard library is
minimal, the language reference is still evolving, and many design choices
remain open to revision.

The claims in this paper are also limited by the absence of user studies.
Without feedback from Limburgian speakers, learners, teachers, or community
organizations, claims about accessibility, cultural impact, authenticity, or
preservation remain speculative. The project can motivate those questions, but
it cannot answer them alone.

Because the language deliberately uses conventional interpreter-language
mechanisms, its contribution to programming-language theory is limited. That is
not a defect in the case study, but it should be stated plainly. The artifact
does not propose a new type system, novel formal semantics, new runtime model,
or optimization technique.

Finally, the cultural framing requires consultation. Regional languages are
living practices, not design assets to be borrowed without accountability.
Before publication, the project should incorporate review by Limburgian
speakers and, where possible, by people involved in Limburgian education,
documentation, or preservation.

% TODO: Add a positionality statement explaining the authors' relationship to
% Limburgian and to the communities discussed in the paper.
